% SHARD-ID: AQM-60-TANGENT-WEYL-OPERATOR-PRINCIPLE
% SHARD-TITLE: Tangent-Weyl Operator Principle and de Rham CSS Prototype
% SHARD-SUMMARY: States the corrected workflow from schemes to tangent-cotangent Weyl operator labels and Hilbert carriers.
% SHARD-SUMMARY: Separates residue-Weyl comparisons and finite \(C^1\)-CSS prototypes from the canonical tangent-Weyl recipe.
% SHARD-SUMMARY: Reworks product-ring, finite-field circle, and formal real-line examples under the operator-selection principle.
% SHARD-KEYWORDS: tangent, cotangent, Weyl, projective representation, stabilizer, de Rham CSS, circle, real line

\section{Tangent-Weyl Operator Principle and de Rham CSS Prototype}
\label{sec:tangent-weyl-operator-principle}

\subsection{Status and Source Anchors}

\begin{sloppypar}
This shard corrects the interpretation of AQM-54--AQM-59.  K\"ahler
differentials are sourced from
\path{references/algebraic_geometry/stacks_project_algebra.tex}, lines
33793--33872 and 34310--34340.  Tangent spaces are sourced from
\path{references/algebraic_geometry/stacks_project_varieties.tex}, lines
2915--3019.  Finite Weyl operators use AQM-28.  The stabilizer/CSS isotropy
dictionary is AQM-08.  The tangent-cotangent carrier convention is
\texttt{CONVENTIONS.md} (ao), with the relative/intrinsic warning in
\texttt{CONVENTIONS.md} (av).  The real CCR comparison is AQM-19.
\end{sloppypar}

The corrected principle is not ``put a Hilbert space on the tangent space'' and
not a separate geometric intermediate step.  The principle is: compute the
pointwise symplectic label spaces, quantise those label spaces by
Weyl-Heisenberg projective unitary representations, and only then ask for a
selection principle choosing commuting Weyl operators.

\subsection{Lamport Dependency Skeleton}

\[
\begin{array}{rcl}
D1&:&k\to A\text{ is fixed},\quad X=\Spec A,\quad P\in X,\quad K=\kappa(P).\\
D2&:&V_P^*=T_{A/k,P}^*X=\Omega^1_{A/k}\otimes_AK.\\
D3&:&V_P=T_{A/k,P}X=\operatorname{Hom}_K(V_P^*,K).\\
D4&:&E_P=V_P\oplus V_P^*.\\
D5&:&\omega((v,\alpha),(w,\beta))=\beta(v)-\alpha(w).\\
D6&:&\rho_P:E_P\to PU(\mathcal H_{\rho_P})
      \text{ is a Weyl-Heisenberg projective representation}.\\
D7&:&e\in E_P\text{ labels a Weyl operator }W_P(e)\text{ on }
      \mathcal H_{\rho_P}.\\
D8&:&E_S=\bigoplus_{P\in S}E_P,\quad
      \rho_S=\bigotimes_{P\in S}\rho_P
      \text{ for finite }S.\\
D9&:&L\subset E_S\text{ isotropic, meaning trivial Weyl commutator}
      \Rightarrow \{W(e):e\in L\}\text{ commutes}.\\
D10&:&\text{A compatible phased-stabilizer character on }L
       \text{ defines a joint eigenspace}.\\
D11&:&C^0=A\to C^1=\Omega^1_{A/k}\to C^2=\Omega^2_{A/k}
      \text{ is a separate finite CSS prototype}.\\
T1&:&D1\text{--}D10\text{ are the tangent-Weyl operator workflow}.\\
T2&:&D11\text{ and residue-Weyl constructions are comparison/prototype layers}.
\end{array}
\]

\subsection{The Local Symplectic Label Space}

Let \(V=V_P=T_{A/k,P}X\) and \(V^*=V_P^*=T^*_{A/k,P}X\), all relative to the
fixed base map \(k\to A\).  Define
\[
  E_P=V\oplus V^*
\]
and
\[
  \omega((v,\alpha),(w,\beta))=\beta(v)-\alpha(w).
\]
The form is bilinear and alternating because
\[
  \omega((v,\alpha),(v,\alpha))=\alpha(v)-\alpha(v)=0.
\]
It is nondegenerate: if \(v\ne0\), choose \(\beta\in V^*\) with
\(\beta(v)\ne0\); if \(\alpha\ne0\), choose \(w\in V\) with
\(\alpha(w)\ne0\).  Pairing with \((0,\beta)\) or \((w,0)\) gives a nonzero
value.  Hence \(E_P\) is the local symplectic label space.

In coordinates \(V=K\,\partial_x\), \(V^*=K\,dx\), this convention gives
\[
  \omega((q,p),(q',p'))=p' q-p q'.
\]
Thus it agrees with the \(dx\wedge dp\) sign used in AQM-58 after that shard's
repair.  It is opposite to the display obtained from the residue-Weyl ordering
\(W(q,p)=T(q)R(p)\); comparison with that ordering requires either reversing
the translation/modulation order or replacing the display momentum by its
negative.

For finite \(K\), the Weyl-Heisenberg step constructs projective unitary
representations of \(E_P\).  With central character
\[
  \psi_K(a)=\exp\!\left(\frac{2\pi i}{p}
    \operatorname{Tr}_{K/\mathbb F_p}(a)\right),
\]
the irreducible finite representation has carrier dimension
\[
  \dim_{\mathbb C}\mathcal H_{\rho_P}=|K|^{\dim_K V}.
\]
This is where Hilbert carriers enter the recipe: they are carriers of the
Weyl-Heisenberg projective representation.  They are not assigned directly to
\(T_{A/k,P}X\).

\subsection{Stabilizer Selection Is the Isotropic Step}

A stabilizer principle chooses label data
\[
  L\subset E_S=\bigoplus_{P\in S}E_P
\]
whose Weyl commutator bicharacter is trivial.  For a common residue field this
is the usual vanishing of the summed symplectic pairing; for mixed finite
residue fields it is the componentwise product of the local trace-character
commutators.  This isotropy condition is exactly what is needed for the
corresponding Weyl operators to commute after the central lifts are fixed.
After the selection principle also fixes a compatible phased-stabilizer
character on \(L\), the selected stabilizer subspace is the joint eigenspace
inside the carrier of \(\rho_S\).

Thus the role previously being described by the wrong geometric word is played
by the stabilizer-selection principle itself.  A maximal isotropic choice is a
Lagrangian-type stabilizer choice, but it is still a choice of commuting Weyl
operators, not an extra canonical step between geometry and representation.

\subsection{Comparison Layers}

The older residue-Weyl comparison attaches to a finite residue field \(K\) the
label shape \(K_q\oplus K_p\) and a carrier such as \(\ell^2(K)\).  This is a
valid comparison layer from AQM-28--AQM-30, but it uses only \(K\).  It does
not see whether \(P\) is smooth, singular, reduced, nonreduced, tangent-rich, or
tangent-zero.

The finite de Rham CSS prototype is also separate.  It uses
\[
  C^0=A\xrightarrow{d_0}C^1=\Omega^1_{A/k}
       \xrightarrow{d_1}C^2=\Omega^2_{A/k}
\]
when these are finite-dimensional \(k\)-vector spaces, chooses a \(C^1\)-basis,
and forms
\[
  R_X=\operatorname{im}d_0,\qquad R_Z=\operatorname{im}d_1^T.
\]
The equation \(d_1d_0=0\) proves CSS isotropy for that auxiliary
chain-complex construction.  It is useful, but it is not the canonical
pointwise tangent-Weyl quantisation of \(X\).

\subsection{Example I: Product Fields}

Let
\[
  A=k^n=k e_1\oplus\cdots\oplus k e_n .
\]
AQM-58 proves that every \(k\)-derivation \(D:A\to M\) vanishes.  Hence
\[
  \Omega^1_{A/k}=0.
\]
At the \(i\)-th point \(P_i\),
\[
  V_{P_i}^*=0,\qquad V_{P_i}=0,\qquad E_{P_i}=0.
\]
The tangent-Weyl representation is therefore trivial at every point: there are
no nontrivial geometric Weyl operator labels to select.  The older residue-Weyl
comparison would attach one \(k\)-level carrier per residue point, but that is
not the corrected tangent-cotangent construction.  The de Rham CSS prototype
also selects no stabilizer because \(C^1=0\).

\subsection{Example II: The Odd Finite-Field Circle}

Let \(k=\mathbb F_q\) with \(q\) odd and
\[
  C=\Spec k[x,y]/(x^2+y^2-1).
\]
For a rational point \(P=(a,b)\), AQM-59 computes
\[
  V_P^*=(k\,dx\oplus k\,dy)/(a\,dx+b\,dy).
\]
The tangent space is
\[
  V_P=\{(u,v)\in k^2:au+bv=0\}.
\]
Since \(a^2+b^2=1\), both spaces are one-dimensional.  Hence
\[
  E_P=V_P\oplus V_P^*
\]
is a two-dimensional symplectic \(k\)-label space, and the finite
Weyl-Heisenberg carrier has dimension \(q\).  For
\[
  C(\mathbb F_3)=\{(1,0),(2,0),(0,1),(0,2)\},
\]
the rational tangent-Weyl support has four local carriers of dimension \(3\).
This is still only kinematics.  A code or state appears only after a principle
selects an isotropic subset of the four-point label space
\[
  \bigoplus_{P\in C(\mathbb F_3)}(V_P\oplus V_P^*).
\]

\subsection{Example III: The Formal Real Affine Line}

Let
\[
  A=\mathbb R[x],\qquad X=\Spec A.
\]
AQM-58 computes
\[
  \Omega^1_{A/\mathbb R}=A\,dx.
\]
At the real point \(P_a=(x-a)\),
\[
  V_{P_a}^*=\mathbb R\,dx_a,\qquad
  V_{P_a}=\mathbb R\,\partial_x|_a .
\]
The local operator-label space is
\[
  E_{P_a}=\mathbb R\,\partial_x|_a\oplus\mathbb R\,dx_a,
\]
with symplectic form
\[
  \omega((r\partial_x,s\,dx),(r'\partial_x,s'\,dx))=s'r-sr'.
\]
The Weyl/CCR representation theory supplies Hilbert carriers for this label
space; the usual \(L^2(\mathbb R)\) realization is an analytic model, not a
stabilizer-selection principle.  It does not choose the free-particle
Hamiltonian, a potential, a domain, boundary conditions, or a distinguished
commuting operator family.

\subsection{Operational Conclusion}

The corrected workflow is:
\[
\begin{array}{ll}
1. & \text{Choose a base map }k\to A\text{ and set }X=\Spec A.\\
2. & \text{Compute points }P\text{ and residue fields }\kappa(P).\\
3. & \text{Compute the relative }T_{A/k,P}X\text{ and }T^*_{A/k,P}X.\\
4. & \text{Build }E_P=T_{A/k,P}X\oplus T^*_{A/k,P}X.\\
5. & \text{Quantise }E_P\text{ by Weyl-Heisenberg projective representations}.\\
6. & \text{Use a separate principle to select isotropic Weyl labels}.\\
7. & \text{Choose compatible stabilizer phases and take the joint eigenspace}.
\end{array}
\]
AQM-61--AQM-63 use this workflow.  AQM-55--AQM-57 remain useful only as
finite \(C^1\)-CSS prototype calculations unless they are explicitly pushed
through the pointwise tangent-Weyl representation layer.
